\section{An absolutely convergent expansion in sites}
\label{sec:expansion}

We now express finite near-critical crossing tests in terms of critical
color laws.  The expansion is in the sites at which a color history is
changed.  Its coefficients can have either sign, and two independent
color replicas may use the same Poisson geometry.  The purpose of this
section is to obtain an absolutely summable bound for all the
coefficients and uniform integrability for each fixed coefficient.
Identification of their limits will then be a separate local problem.

Recall the fixed rectangle $Q_0$ and the normalization from
Lemma~\ref{lem:pivotal-scale}:
\begin{equation}
 q_e=\frac1{\EE N_e(Q_0)},\qquad
 a_e=A_e(b_e,1),\qquad m_eq_ea_e\asymp1,
 \qquad q_e\longrightarrow0.
 \label{eq:auxiliary-scale}
\end{equation}
Fix finitely many polygonal quads, one or two replica labels
$r\in\{1,\ldots,R\}$, and, for each label, finitely many real times.
Let $F$ be any function of the corresponding crossing indicators with
$|F|\leq1$.  In this section the density at time $t$ is
$1/2+tq_e$.  For times in a fixed bounded interval these densities lie
in $[0,1]$ for all sufficiently small $e$.

\subsection{Exact color histories and signed kernels}

For a single replica with time vector
$\mathbf t=(t_1,\ldots,t_n)$, a history is an element of
$\mathcal H=\{0,1\}^n$.  Write $\mathbf0$ and $\mathbf1$ for its
two constant elements and put
$P_0=(\delta_{\mathbf0}+\delta_{\mathbf1})/2$.
If $T\geq\max_i|t_i|$, define the signed measure
\begin{equation}
 w_{\mathbf t}
 =\int_{-T}^{T}\delta_{(\mathbf1_{\{s\leq t_i\}})_{i=1}^n}\,ds
       -T\delta_{\mathbf0}-T\delta_{\mathbf1}.
 \label{eq:history-weight}
\end{equation}
This measure is independent of the chosen $T$ once $T$ is large enough.
It has total mass zero and total variation at most $4T$.  The law of
$(\mathbf1_{\{U\leq1/2+t_iq_e\}})_{i=1}^n$, for a uniform mark $U$,
is exactly
\begin{equation}
                         P_0+q_ew_{\mathbf t}.
 \label{eq:exact-history-law}
\end{equation}
For example, in increasing time order the masses of
$w_{\mathbf t}$ at the all-one and all-zero histories are respectively
$\min_i t_i$ and $-\max_i t_i$; its other nonzero masses are successive
time differences.  Formula~\eqref{eq:history-weight} also handles tied
times without a choice of ordering.  It shows that
\begin{equation}
 \|w_{\mathbf t}-w_{\mathbf t'}\|_{\TV}
       \leq2\sum_i|t_i-t_i'|.
 \label{eq:history-continuity}
\end{equation}
Thus the finite labeled history space can be kept fixed while the times
vary.  Different replicas may have different time lists.

Choose a fixed box $B$ containing every tested quad with a positive
margin.  In the Poisson model choose a geometry event $\Gamma_e$ which
ensures that all cells meeting the tested region have their nuclei in
$B$ and that their relevant portions are determined by the nuclei in
$B$.  A finite covering by occupancy squares inside $B$ gives such an
event, increasing under addition of nuclei, with
\begin{equation}
                  \PP(\Gamma_e^c)\leq C\exp(-c/e^2).
 \label{eq:expansion-cutoff}
\end{equation}
We multiply the test by $\mathbf1_{\Gamma_e}$; we do not condition on
$\Gamma_e$.  Off this event its value is defined to be zero, so it is
a function of the colors of the finitely many sites in $B$.  In the
triangular model the fixed margin makes this localization deterministic
for all sufficiently small meshes, and we take $\Gamma_e$ to be the
whole sample space.

The counting volume, which converges to planar area, is
\begin{equation}
 \sigma_e(dx)=
 \begin{cases}
 dx,&\text{in the Poisson model},\\
 m_e^{-1}\displaystyle\sum_{x\in T_e}\delta_x(dx),
       &\text{in the triangular model}.
 \end{cases}
 \label{eq:counting-volume}
\end{equation}
For an ordered list $\mathbf x=(x_1,\ldots,x_k)$ of distinct sites in
$B$, choose nonempty sets $I_i\subseteq\{1,\ldots,R\}$.
For every $r\in I_i$ let $w_{i,r}$ be a signed measure of total mass
zero on that replica's history space.  Define
\begin{equation}
 K_e(F;\mathbf x,\mathbf I,\mathbf w)
 =\EE\left[
   \int \mathbf1_{\Gamma_e}
        F\bigl(\omega[\mathbf h]\bigr)
        \prod_{i=1}^k\prod_{r\in I_i}dw_{i,r}(h_{i,r})
       \right].
 \label{eq:site-kernel}
\end{equation}
Here, in the Poisson case, $x_1,\ldots,x_k$ are first added as nuclei
to an ordinary Poisson process; in the triangular case the tiling is
unchanged.  All unselected site--replica coordinates have independent
fair, constant histories.  At the selected coordinates
$\omega[\mathbf h]$ uses the histories being integrated.  The
expectation averages the remaining fair bits and, in the Poisson case,
the Poisson process.  An unselected replica at an inserted nucleus is
therefore still fair.  Insertion of a nucleus in this formula is the
Mecke representation of a sum over existing sites; the signed operation
itself changes colors and does not define geometric pivotality.

Write $B^k_{\ne}$ for the set of distinct tuples, and
$d(\mathbf I)=\sum_i|I_i|$.  Multiplying
\eqref{eq:exact-history-law} over sites and replica labels, and then
using the Poisson factorial moment formula, gives
\begin{equation}
 \EE[\mathbf1_{\Gamma_e}F]
 =\EE_0[\mathbf1_{\Gamma_e}F]
  +\sum_{k\geq1}\frac{m_e^k}{k!}
    \sum_{\substack{\varnothing\ne I_i\subseteq\{1,\ldots,R\}\\1\leq i\leq k}}
    q_e^{d(\mathbf I)}
    \int_{B^k_{\ne}}K_e(F;\mathbf x,\mathbf I,\mathbf w)
                                      \,d\sigma_e^{\otimes k}(\mathbf x).
 \label{eq:site-expansion}
\end{equation}
In this identity $w_{i,r}=w_{\mathbf t^{(r)}}$, and $\EE_0$ uses fair
constant histories everywhere.  The left side uses the actual
near-critical histories.  For each fixed $e$ the lattice expansion is
finite.  The Poisson identity and its initial interchange of expectation
and summation follow from the exponential moments of the number of
nuclei in $B$.  The stronger, mesh-uniform summability needed for taking
limits is the content of the next proposition.

\begin{proposition}[Uniform bound for the site expansion]
\label{prop:expansion-bound}
Fix a box $B$, finitely many polygonal crossing tests, and $R\leq2$
replica labels.  Let $g\in(1,2)$ be an exponent in the lower four-arm
bound of Lemma~\ref{lem:arm-inputs}.  There is a constant $C$, depending
only on this data, such that, for every $k\geq1$, every choice of
nonempty $I_i$, and centered weights $w_{i,r}$,
\begin{equation}
 \frac{(m_eq_e)^k}{k!}
   \int_{B^k_{\ne}}
       |K_e(F;\mathbf x,\mathbf I,\mathbf w)|
                    \,d\sigma_e^{\otimes k}(\mathbf x)
 \leq C^k k^{-(1-g/2)k}
             \prod_{i,r\in I_i}\|w_{i,r}\|_{\TV},
 \label{eq:coefficient-bound}
\end{equation}
uniformly for all sufficiently small meshes and all $|F|\leq1$ of the
specified crossing indicators.  Consequently the absolute sum in
\eqref{eq:site-expansion} is bounded by
$\sum_{k\geq1}C_T^k k^{-(1-g/2)k}<\infty$ when all times have absolute
value at most $T$.  For two replicas, the total absolute contribution
of terms with $|I_i|=2$ for at least one $i$ is $O_T(q_e)$.
\end{proposition}

\begin{proof}
We first establish a pointwise geometric bound and then integrate it.
Put $W(\mathbf w)=\prod_{i,r\in I_i}\|w_{i,r}\|_{\TV}$; the assertion
is immediate if one of these factors is zero.
For fixed geometry and unselected fair bits, the signed integral in
\eqref{eq:site-kernel} vanishes if $F$ is independent of any one
selected history, because its weight has mass zero.  Thus a necessary
condition for a nonzero integral is the following: at every selected
coordinate there are settings of the inserted histories for which
changing that coordinate changes at least one tested crossing.
The settings witnessing this condition may differ from coordinate to
coordinate.  We use only the resulting arm tests in regions containing
no inserted site, where all the colors are the original fair bits.
In particular, no enumeration of all global history assignments is
required.

Our geometric objective is one four-arm factor $a_e$ for each inserted
site, with the remaining cost expressed as a product of integrable
singularities in intersite distances and distances to polygon vertices
and marks.  The
merging forest will select the distances and the disjoint annuli that
supply these factors.

\paragraph{The merging forest.}
Organizing nearby sites into merging clusters and assigning arm events
to the annuli between successive scales has a methodological
predecessor in the spectral-sample analysis of Garban, Pete, and
Schramm~\cite[Sections~4.2--4.3]{gps-fourier}.  Here we need estimates
for signed history kernels rather than squared Fourier coefficients,
and give the required argument.
Place fixed roots at all polygon vertices and marks, identifying
coincident roots, and denote their number by $J$.  Identify these roots
with a common sink for the purpose of Kruskal's algorithm.  On the
remaining graph use the ordinary Euclidean distances between all
insertion sites and roots, clipped below at $b=b_e$.  Starting at scale
$b$, accept an edge when its endpoints are in distinct components, and
continue through a fixed scale $L<1$.  Ties are resolved by a fixed
order.  The constant $L$ is chosen small enough, in terms of the fixed
polygons, that balls of radius $O(L)$ do not meet two nonincident sides
of any one polygon.  The constants in the probes below will also be
chosen in terms of the incident angles.  Record the ordinary
accepted edges, with clipped lengths $\ell_f$; the virtual
identification of the roots is not a recorded edge.

A component containing the sink is called grounded.  A \emph{life} is
a maximal interval $[u,v]$ during which an ungrounded component has the
same set of insertion sites.  Each merge ends the participating
ungrounded lives and, if the merged component is ungrounded, starts a
new life.  A remaining life ends at $L$.  Zero-length lives caused by
ties can be discarded.  There are at most $2k$ lives and at most $k$
recorded edges: every accepted edge decreases the number of
ungrounded components by one.  The ordinary recorded graph is a forest
and each of its components contains at most one root, since an edge
joining two rooted components would already make a cycle through the
sink.

For a life $C$ on $[u,v]$, retain one representative site $x_C$ and
write $D_C$ for the diameter of its site set.  Take a buffered annular
probe centered at $x_C$, with radii
\begin{equation}
                C_0\max(u,D_C)\quad\hbox{and}\quad c_0v,
 \label{eq:life-probe}
\end{equation}
where $C_0$ is sufficiently large and $c_0$ sufficiently small.
If these radii are reversed, impose no probe and assign it cost one.
We use the convention $A_e(r,R)=1$ when $r\geq R$.

The necessary influence of the representative site gives an arm test
on every nonempty probe.  The changing cell is localized inside its
inner boundary, and the other inserted histories do not enter the
probe, on the occupancy events of Lemma~\ref{lem:localization}.
For each tested quad separately the probe meets no side or at most one
straight side: all vertices and marks are outside the outer buffer,
and the constants are chosen in terms of the finitely many polygon
angles.  When a side is encountered partway through the probe, use
four arms up to the distance to that side and three half-plane arms
thereafter.  Lemma~\ref{lem:pivotal-arms} and the straight-side estimate
bound this cost by
\begin{equation}
                  C_1 A_e(C_0\max(u,D_C),c_0v).
 \label{eq:life-arm-cost}
\end{equation}
A purported influence outside the quad is either impossible after
localization or has the same straight-side test.  Different polygons
may overlap; the argument just chooses one quad and one replica for
the influence at this probe.  There are only a fixed number of such
choices per life.  Even if its witness uses different times, the bits
on the probe are critical and constant in time.

The probes can be chosen with disjoint buffers.  If a life $D$ descends
from a life $C$, then $x_D$ belongs to the site set of $C$ and
$v_D\leq u_C$.  Its outer buffer is thus within distance
$D_C+O(c_0u_C)$ of $x_C$, strictly inside the inner boundary of $C$'s
probe when $C_0$ is large.  For two components with
incomparable lives, let $v$ be the larger of their death scales.  Their
site sets are at distance at least $v$: otherwise Kruskal would process
an edge between them before scale $v$, forcing the component with that
death scale either to contain a site of the other component already at
its birth or to change its site set before its death.  If the other
component has become grounded, this edge would instead ground it
before its death.  Each alternative contradicts the definition of the
life.  Their outer radii are at most $c_0v$, so their buffers are
disjoint after decreasing $c_0$.  This also applies when the earlier
component has become grounded.  The argument concerns the entire
component sets, so a change of representative causes no loss.

In Poisson percolation each localized probe uses only the Poisson
points and bits in its buffer.  Forced nuclei all lie outside these
buffers.  If an occupancy test fails, use the outermost-failure
decomposition of Lemma~\ref{lem:localization}: omit the failed inner
part and retain the arm probe beyond it, with disjoint margins.  Its
alternatives remain measurable in the original probe buffer, and
their sum has the same upper bound~\eqref{eq:life-arm-cost}.  This also
accounts for a cell reaching far from its nucleus.  Thus the probe
costs multiply.  On the lattice their site sets are disjoint when the
constant in $b_e$ is sufficiently large.  For common-geometry replicas
only one replica is used in each probe; disjoint spatial buffers still
give the required independence.  We obtain
\begin{equation}
 |K_e(F;\mathbf x,\mathbf I,\mathbf w)|
 \leq C_2^k W(\mathbf w)
       \prod_{C\text{ a life}}
          A_e(C_0\max(u_C,D_C),c_0v_C).
 \label{eq:probe-product}
\end{equation}

\paragraph{Diameter losses and telescoping.}
The component diameters in~\eqref{eq:probe-product} must be removed
without a factor growing faster than exponentially in $k$.
Quasi-multiplicativity and the power lower bound show that replacing
the factor for a life by $A_e(u_C,v_C)$ costs at most a fixed constant
times
\[
 \exp\!\left(g\int_{u_C}^{v_C}
                         \mathbf1_{\{t<D_C\}}\frac{dt}{t}\right).
\]
The fixed factors $C_0,c_0$ cost another constant per life, including
the case of an empty probe.  Each ungrounded component is joined by
its internal recorded edges, so
$D_C\leq\sum_{f\text{ internal to }C}\ell_f$.
At a time $u_C<t<v_C$ those edges all have $\ell_f\leq t$.
The currently live ungrounded components are disjoint; hence an edge
is charged by at most one life at any time.  Consequently
\begin{align}
 \sum_C\int_{u_C}^{v_C}\mathbf1_{\{t<D_C\}}\frac{dt}{t}
 &\leq\sum_C\sum_{f\text{ internal to }C}
                         \int_{u_C}^{v_C}\frac{\ell_f}{t^2}\,dt
 \leq\sum_f\int_{\ell_f}^{\infty}\frac{\ell_f}{t^2}\,dt
 \leq k.
 \label{eq:diameter-charge}
\end{align}
Clipping edges below $b$ preserves the diameter inequality, so this
estimate includes clusters of arbitrarily close insertions.

Set $p_e(t)=A_e(t,L)$.  Up to a constant for each life,
$A_e(u,v)=p_e(u)/p_e(v)$.  Initially there are $k$ factors $p_e(b)$.
When two ungrounded components merge at scale $\ell$, two denominators
$p_e(\ell)$ are canceled by one numerator, leaving one inverse
factor.  When an ungrounded component joins the sink there is again one
inverse factor.  Lives surviving to $L$ end with $p_e(L)=1$.
There are therefore no extra terminal factors, and
\begin{align}
 \prod_C A_e(u_C,v_C)
 &\leq C_3^k A_e(b,L)^k\prod_f A_e(\ell_f,L)^{-1}
 \leq C_4^k a_e^k\prod_f\ell_f^{-g}.
 \label{eq:forest-bound}
\end{align}
The comparison of $A_e(b,L)$ with $a_e$ uses the fixed positive scale
$L$.  Combining~\eqref{eq:probe-product}--\eqref{eq:forest-bound}
proves the required pointwise bound.  No estimate is needed for the
color law inside a cluster of nearby insertions.

\paragraph{Integrating the forest.}
Put $s=k^{-1/2}$.  Call a recorded edge short when $\ell_f<s$.
There are at most $k$ recorded edges.  If there are $j$ short edges,
the remaining factors in~\eqref{eq:forest-bound} are at most
$k^{g(k-j)/2}$; this remains an upper bound when the forest has fewer
than $k$ edges.  The short edges form a forest with at most one root in
each component.  We may overcount their choices by
$\binom{(k+J)^2}{j}$, ignoring all conditions imposed by Kruskal's
algorithm.

Uniformly in a point $y$ of the plane, the one-edge integral satisfies
\begin{equation}
 \int_B (|x-y|\vee b)^{-g}
           \mathbf1_{\{|x-y|\vee b<s\}}\,d\sigma_e(x)
                  \leq C_5s^{2-g}.
 \label{eq:short-edge-integral}
\end{equation}
For area measure this follows by polar integration.  For lattice
volume, $\sigma_e(B(y,r))\leq C(r^2+e^2)$, and $b\geq ce$ permits a
dyadic-shell summation from $b$ to $s$ with the same bound.  This is
uniform also when $y$ is a fixed root off the lattice.  If $b\geq s$
the integral is zero.  Allowing equal lattice sites only enlarges this
bound; the atom at $x=y$, when present, has mass
$O(e^2)b^{-g}=O(b^{2-g})$.

Integrate the leaves of every short-edge tree.  Each edge gives the
factor in~\eqref{eq:short-edge-integral}; a tree without a fixed root
leaves one free center, whose integral is at most a constant depending
on $B$.  A rooted tree leaves no such integration.  We conclude that
\begin{align}
 \frac1{k!}\int_{B^k_{\ne}}\prod_f\ell_f^{-g}
                              \,d\sigma_e^{\otimes k}
 &\leq \frac{C_6^k}{k!}
       \sum_{j=0}^k k^{g(k-j)/2}
          \binom{(k+J)^2}{j}s^{(2-g)j} \notag\\
 &=\frac{C_6^k k^{gk/2}}{k!}
       \sum_{j=0}^k\binom{(k+J)^2}{j}k^{-j}
 \leq C_7^k k^{-(1-g/2)k}.
 \label{eq:factorial-forest}
\end{align}
For the last inequality, extend the sum to the binomial expansion of
$(1+1/k)^{(k+J)^2}\leq C_J^k$ and use $k!\geq(k/\mathrm e)^k$.
Thus both the number of possible forests and the small-scale
singularities have been included in the factorial bound.

Finally $(m_eq_e)a_e$ is bounded by~\eqref{eq:auxiliary-scale}, giving
\eqref{eq:coefficient-bound}.  For $R\leq2$ there are at most $3^k$
label assignments.  On a bounded time range all variation norms are
bounded, and $q_e^{d(\mathbf I)-k}\leq1$ once $q_e\leq1$.
The asserted summable majorant follows.  If some $|I_i|=2$, this last
factor is at most $q_e$.  The same summable majorant, multiplied by
$q_e$, bounds the entire repeated-label contribution.
\end{proof}

\subsection{Location cutoffs and fixed coefficients}

The coefficient estimate sums over all locations, including clusters
and positions near polygon edges.  For the local comparison in the
next section we need to remove such positions at arbitrarily small
cost.  A slightly stronger integrability estimate provides this.

\begin{lemma}[Uniform integrability in the insertion locations]
\label{lem:location-ui}
Fix $k$, the polygonal test family, $R\leq2$, and an upper bound on
all the variation norms in~\eqref{eq:site-kernel}.  There are $p>1$ and
$C_k<\infty$ such that, after extending the kernel by zero to coincident
tuples,
\begin{equation}
 \int_{B^k}
   \left|(m_eq_e)^kK_e(F;\mathbf x,\mathbf I,\mathbf w)\right|^p
                      \,d\sigma_e^{\otimes k}(\mathbf x)
                   \leq C_k.
 \label{eq:kernel-lp}
\end{equation}
Let $Z$ be the union of all tested polygon boundaries and $\partial B$.
For $\rho>0$ let $\mathcal B_\rho\subset B^k$ consist of tuples for
which $|x_i-x_j|<\rho$ for some $i\ne j$ or
$\dist(x_i,Z)<\rho$ for some $i$.  Then
\begin{equation}
 \lim_{\rho\downarrow0}\limsup_{e\downarrow0}
    (m_eq_e)^k\int_{\mathcal B_\rho}
       |K_e(F;\mathbf x,\mathbf I,\mathbf w)|
                       \,d\sigma_e^{\otimes k}(\mathbf x)=0.
 \label{eq:location-cutoff}
\end{equation}
The conclusions are uniform in the indicated tests and weights.
They remain true after omitting the geometry cutoff $\Gamma_e$ from
the definition of the kernel.
\end{lemma}

\begin{proof}
Choose $p>1$ with $gp<2$.  For fixed $k$, there are only finitely many
ordinary forests on the $k$ sites and the $J$ fixed roots.  The
pointwise estimate~\eqref{eq:forest-bound}, multiplied by
$(m_eq_e)^k$, is bounded by a constant times the sum over all such
forests having at most one root per component of
$\prod_f(|x_{i(f)}-x_{j(f)}|\vee b)^{-g}$, with a root coordinate
interpreted as fixed.  Raising this finite sum to the power $p$ costs
only a constant depending on $k$.  Leaf integration, now with exponent
$gp<2$ and with distances bounded by a fixed constant containing $B$
and all roots, bounds each forest integral uniformly in $e$.
The lattice estimate is the same dyadic calculation as
\eqref{eq:short-edge-integral}.  This proves~\eqref{eq:kernel-lp}.

H\"older's inequality bounds the expression before the limits in
\eqref{eq:location-cutoff} by
$C_k^{1/p}\sigma_e^{\otimes k}(\mathcal B_\rho)^{1-1/p}$.
For area measure this volume tends to zero as $\rho\downarrow0$.
For lattice volume its mesh limit is the same: the finitely many
polygon edges and the pairwise diagonals have zero area in the relevant
product space, and their $\rho$-neighborhoods can be bounded by finite
unions of boxes with arbitrarily small total volume as $\rho$ decreases.
This proves the ordered-limit assertion and includes vertices, marks,
and box edges.

It remains to justify removal of $\Gamma_e$.  Adding a fixed number
of nuclei can only improve its occupancy requirements.  Thus the
difference between truncated and untruncated kernels is bounded
uniformly in the locations by
$C_k\exp(-c/e^2)$.  The normalization $(m_eq_e)^k$ grows only
polynomially in $1/e$, by the lower arm bound in
Lemma~\ref{lem:arm-inputs}.  The normalized difference therefore
converges uniformly to zero, also in $L^p$.
\end{proof}

For later calibration, consider a single time and the centered weight
$\Delta=\delta_1-\delta_0$, enlarging $B$ to contain $Q_0$ with a
margin if necessary.  For the increasing crossing test of $Q_0$
the corresponding one-site kernel is its nonnegative color-pivotal
probability.  Russo's identity and the definition of $q_e$ give
\begin{equation}
 (m_eq_e)\int_B
      K_e(\mathbf1_{C_{Q_0}};x,\Delta)\,d\sigma_e(x)=1+o(1).
 \label{eq:first-coefficient-one}
\end{equation}
Here the single replica label is suppressed.  Without the box cutoff,
the integral over all sites is exactly $q_e\EE N_e(Q_0)=1$.
For the displayed truncated version its difference is negligible:
localization gives exponentially small failure probability and the
number of cells meeting a fixed compact set has moments growing at
most polynomially in $1/e$, as in Lemma~\ref{lem:pivotal-scale}.
Cauchy--Schwarz then also controls the lost pivotal count.

We have reduced all finite-time tests to an absolutely convergent
series of critical signed kernels.  Terms using both replicas at a
single site disappear, and each remaining fixed-degree integral may
be restricted to compact sets of separated locations away from the
polygon boundaries.  These are precisely the conditions under which
the local comparison in the next section will identify its limit.
